\documentclass[11pt,a4paper]{article}
\usepackage{curso}
\begin{document}
\portada{Ejercicios --- Lección}{0.3}{Repaso de probabilidad}

\section{Ejercicios básicos}
\begin{enumerate}[label=\textbf{B\arabic*.}]
\item (Pensum, Ap. A.1) Calcule la media y la covarianza de $x=Ay+b$ cuando $y\sim\N(0,I)$.
\item (Pensum, Ap. A.5) Calcule la entropía de una Bernoulli($p$) y de una gaussiana $d$-dimensional $\N(\mu,\Sigma)$.
\item Calcule $\KL(\N(0,1)\|\N(0,4))$ y $\KL(\N(0,4)\|\N(0,1))$.
\item ¿Cuántos lanzamientos se necesitan para estimar $p\approx0{,}5$ con error estándar $0{,}01$?
\end{enumerate}

\section{Ejercicios intermedios}
\begin{enumerate}[label=\textbf{I\arabic*.}]
\item Sea $s\in\{-1,+1\}$ con $P(s=+1)=p$. Exprese $\E[s]$ y $\Var[s]$ en función de $p$. Si $p=\sigma(2u)$ con $\sigma(z)=1/(1+e^{-z})$, muestre que $\E[s]=\tanh u$.
\item Dada $\Sigma=\begin{pmatrix}2&1\\1&2\end{pmatrix}$, encuentre una $A$ tal que $AA^\top=\Sigma$ (Cholesky) y otra distinta (raíz simétrica). ¿Generan la misma distribución?
\item Demuestre $\KL(p\|q)\ge0$ usando $\ln z\le z-1$.
\item Calcule la entropía de una variable categórica uniforme con $K$ valores y la de un dado con probabilidades $(0{,}1;0{,}25;0{,}15;0{,}2;0{,}1;0{,}2)$.
\end{enumerate}

\section{Ejercicios avanzados}
\begin{enumerate}[label=\textbf{A\arabic*.}]
\item Derive la KL entre gaussianas $d$-dimensionales: $\KL(\N_0\|\N_1)=\frac12\big[\Tr(\Sigma_1^{-1}\Sigma_0)+(\mu_1-\mu_0)^\top\Sigma_1^{-1}(\mu_1-\mu_0)-d+\ln\frac{\det\Sigma_1}{\det\Sigma_0}\big]$.
\item Muestre que entre todas las densidades en $\R$ con varianza $\sigma^2$, la gaussiana maximiza la entropía diferencial (multiplicadores de Lagrange).
\item El estado estacionario del OU es $\N(A^{-1}b,\beta^{-1}A^{-1})$. Si se promedian $n$ muestras independientes para estimar $x^*=A^{-1}b$, ¿cuál es la covarianza del estimador? ¿Cómo escala el error con $\beta$ y $n$?
\end{enumerate}

\section{Soluciones}
\begin{solucion}
\textbf{B1.} $\E[x]=b$, $\mathrm{Cov}(x)=AA^\top$ (véanse las notas).

\textbf{B2.} $h(p)=-p\log_2p-(1-p)\log_2(1-p)$ bits. $h(\N(\mu,\Sigma))=\frac12\ln\det(2\pi e\Sigma)$ nats.

\textbf{B3.} $\ln2+\frac18-\frac12=0{,}318$ nats; $\ln\frac12+\frac42-\frac12=0{,}807$ nats.

\textbf{B4.} $\sqrt{0{,}25/n}=0{,}01\Rightarrow n=2500$.

\textbf{I1.} $\E[s]=2p-1$; $\Var[s]=1-(2p-1)^2=4p(1-p)$. Con $p=\sigma(2u)$: $2\sigma(2u)-1=\frac{1-e^{-2u}}{1+e^{-2u}}=\tanh u$.

\textbf{I2.} Cholesky: $L=\begin{pmatrix}\sqrt2&0\\1/\sqrt2&\sqrt{3/2}\end{pmatrix}$. Raíz simétrica: con valores propios $3$ y $1$ y vectores $(1,1)/\sqrt2$, $(1,-1)/\sqrt2$: $\Sigma^{1/2}=\frac12\begin{pmatrix}\sqrt3+1&\sqrt3-1\\\sqrt3-1&\sqrt3+1\end{pmatrix}$. Ambas dan $AA^\top=\Sigma$, luego la misma distribución (difieren en una rotación $A'=AQ$).

\textbf{I3.} $-\KL(p\|q)=\E_p[\ln(q/p)]\le\E_p[q/p-1]=\int q-\int p=0$.

\textbf{I4.} Uniforme: $\log_2K$. Dado: $-\sum p\log_2p=2(0{,}1\cdot3{,}322)+0{,}25\cdot2+0{,}15\cdot2{,}737+2(0{,}2\cdot2{,}322)=0{,}664+0{,}5+0{,}411+0{,}929\approx2{,}50$ bits (máximo $\log_26=2{,}585$).

\textbf{A1.} Use $\E_0[(x-\mu_1)^\top\Sigma_1^{-1}(x-\mu_1)]=\Tr(\Sigma_1^{-1}\Sigma_0)+(\mu_0-\mu_1)^\top\Sigma_1^{-1}(\mu_0-\mu_1)$ y $\E_0[(x-\mu_0)^\top\Sigma_0^{-1}(x-\mu_0)]=d$, junto con los logaritmos de las normalizaciones.

\textbf{A2.} Maximizar $-\int p\ln p$ con $\int p=1$, $\int x p=\mu$, $\int(x-\mu)^2p=\sigma^2$ da $\ln p=\lambda_0+\lambda_1x+\lambda_2(x-\mu)^2$, es decir una gaussiana; la concavidad de la entropía garantiza máximo.

\textbf{A3.} $\mathrm{Cov}(\bar x)=\frac{1}{n\beta}A^{-1}$. El error cuadrático medio total es $\Tr(A^{-1})/(n\beta)$: baja con más muestras y con menos temperatura (más $\beta$). Si las muestras están correlacionadas en el tiempo, $n$ se reemplaza por el tamaño efectivo de muestra (lección 2.6).
\end{solucion}

\section{Ejercicios adicionales (sin solución)}
\begin{enumerate}[label=\textbf{P\arabic*.}]
\item Genere $10^6$ muestras de Bernoulli con \texttt{jax.random} y estime su entropía (laboratorio L1).
\item Verifique numéricamente la fórmula de KL entre gaussianas por Monte Carlo.
\item Grafique $\KL(\N(0,1)\|\N(m,s^2))$ y $\KL(\N(m,s^2)\|\N(0,1))$ en el plano $(m,s)$. ¿Dónde difieren más?
\end{enumerate}
\end{document}
